\documentclass[10pt,a4paper]{amsart}
\usepackage[margin=19mm]{geometry}
\usepackage{amsmath,amssymb,mathtools}
\usepackage[scr=rsfs,cal=euler]{mathalfa}
\usepackage{xfrac}
\usepackage[unicode,hidelinks,bookmarks=false]{hyperref}
\newcommand{\ie}{i.e.\ }
\newcommand{\cf}{cf.\ }
\newcommand{\resp}{respectively\ }
\newcommand{\Subsection}{\subsection}
\providecommand{\nobreakdash}{\nobreak\hbox{-}}
\setlength{\emergencystretch}{2em}
\multlinegap0pt
\usepackage{bm}
\usepackage{bbm}
\usepackage{dsfont}
\usepackage{xspace}
\usepackage{cancel}
\usepackage{mathtools}
\usepackage{amsthm}
\makeatletter
\@ifundefined{thm}{\newtheorem{thm}{Thm}}{}
\@ifundefined{lem}{\newtheorem{lem}[thm]{Lemma}}{}
\makeatother
\def\upd{\mathrm{d}}
\newcommand{\R}{{\mathbb{R}}}
\newcommand{\dd}{\mathrm{d}}
\title[Stochastic Homogenization of Non-local Hamilton-Jacobi-Bellm]{Stochastic Homogenization of Non-local Hamilton-Jacobi-Bellman equations}
\author{Wenjia Jing, Qi Zhang}
\date{Source: arXiv:2608.06814v1 (CC BY 4.0)}
\begin{document}
\maketitle

\noindent\emph{Source and licence.} Stochastic Homogenization of Non-local Hamilton-Jacobi-Bellman equations. Version arXiv:2608.06814v1 (CC BY 4.0), distributed under the Creative Commons Attribution 4.0 International licence, as recorded in the arXiv licence metadata for this exact version and shown on the abstract page https://arxiv.org/abs/2608.06814v1. Full rights notice: licenses/arxiv-2608.06814-CC-BY-4.0.txt.\par\smallskip

\noindent\textit{Editorial scope.} Counted unit: the Lem whose source label is \texttt{Ecv} in the article (source0.tex lines 1208--1214, source label \texttt{Ecv}), delivered together with the article's own complete original proof of it (source0.tex lines 1215--1306, one \texttt{proof} environment of 92 lines). No mathematics is written, repaired or re-derived: statement and proof are the article's own text line for line. Required context retained uncounted: eq:cRdef (lines 344--347); eq:sublifted (lines 1075--1078); Gardv (lines 1174--1180). Every definition, display and label of those lines is retained unchanged. Omitted from this package and not redistributed: 1--343, 348--1074, 1079--1173, 1181--1207, 1307--1867. Nothing inside a retained excerpt is omitted or altered: every retained body is delivered exactly as the article's lines are written, so no omission receipt of any kind accompanies this package. Citations: the retained excerpts carry no citation command at all, so no bibliography entry of the article is transcribed and no reference is redistributed. Reference adaptation: none was needed and none was made; every reference inside the delivered unit resolves inside it, so no unresolved reference or citation is produced. Printed slips kept literal: no printed word, symbol, hypothesis or number of the counted statement or of its proof was altered. Layout and class: the article's own class is \texttt{amsart} with packages including fullpage, amsmath, amsthm, verbatim, amssymb, amsfonts, amscd, graphicx, bbm, graphics, mathtools, mathrsfs, esint, xcolor; the wrapper is the lane's pinned \texttt{amsart} editorial wrapper at 10pt on A4, which restates the theorem-like environments the retained text needs and adds the article's own macro definitions that the retained text needs (\texttt{\textbackslash R}, \texttt{\textbackslash dd}, \texttt{\textbackslash upd}), with extra wrapper packages (bm, bbm, dsfont, xspace, cancel, mathtools). The delivered numbering is the unit's own. Assets: no retained span contains a figure environment, a graphics-inclusion command, a PDF-page inclusion, a table or a TikZ picture; no image file is copied into this package, the package declares no asset dependency and no image file is redistributed with it. Licence: version arXiv:2608.06814v1 of this article is distributed under the Creative Commons Attribution 4.0 International licence, as recorded in the arXiv licence metadata for this exact version and shown on the abstract page https://arxiv.org/abs/2608.06814v1. A full rights notice is delivered at licenses/arxiv-2608.06814-CC-BY-4.0.txt. Characters: every retained excerpt is pure ASCII, so the delivered engine, XeLaTeX with its default Latin Modern text font, reports no missing character.
\begin{equation}
\label{eq:cRdef}
    \mathcal{R}g(x) = \int_{\R^n} R(x-y)g(y)\,\,\mathrm{d}y, \qquad R(x) = \frac{\Theta(\hat{x})}{|x|^{n+\alpha-2}},
\end{equation}

\begin{equation}
\label{eq:sublifted}
    \mathfrak{J}v(\omega) + H(p + v(\omega), \omega) \leq \overline{H}(p),
\end{equation}

\begin{lem}\label{Gardv}
The sequence $\{v_{k}\}$ is weakly compact in $L^{\beta}(\Omega;\R^n)$, and any weak limit $v = (v^1,\dots,v^n) \in L^{\beta}(\Omega;\R^n)$ is a zero mean and curl-free vector field in the following sense,
$$
\mathbf{E}[v] = 0, \quad \mathfrak{D}\times v = \left(\mathfrak{D}_j v^i - \mathfrak{D}_i v^j\right)_{i,j} = 0.
$$
and $v$ satisfies \eqref{eq:sublifted} almost surely. 
\end{lem}

\begin{lem}\label{Ecv} Given $p\in \R^n$, let $v$ be the $L^\beta(\Omega)$ random vector determined in Lemma \ref{Gardv}, and let $\widetilde{v}$ be the random field $\widetilde{v}(x,\omega)=v(\tau_x \omega)$. Then $\widetilde{v}$ is in $L^\beta_{\rm loc}(\R^n)$ uniformly in $\Omega$ in the following sense: for any $R\ge 1$ there is a constant $C_{R} < \infty$, which depends only on $R,p,c_1,c_2,\beta$, such that 
\begin{equation}
\label{eq:unifbddv}
     \sup_{x\in\mathbb{R}^n}\int_{B_R(x)} |\widetilde v(y,\omega)|^{\beta} \,\mathrm{d}y  \le C_{R,p}
\end{equation}
holds $\mathbf{P}$-almost surely.
\end{lem}

\begin{proof} Let $\phi(x) = c(\sqrt{1+|x|^2})^{-n-\delta}$ for some $\delta>0$ to be fixed later and $c>0$ is chosen so that $\int_{\R^n} \phi =1$. In view of \eqref{eq:sublifted} and assumption $(\bf{A})$, we get
\begin{equation}
    \int_{\mathbb{R}^n}(\mathscr{J}\widetilde v (x,\omega) + c_{1}|p + \widetilde v(x,\omega)|^{\beta} - c_{2})\phi(x)\,\mathrm{d}x \leq \overline{H}(p).
\end{equation}
Recall that $\mathscr{J}\widetilde v=\nabla \cdot (\mathcal{R}\widetilde v)$ and that $\mathcal{R}$ is a regular integral operator; see \eqref{eq:cRdef}. Using integration by parts and some basic inequalities, we get
\begin{equation}\label{vLpE}
    \int_{\mathbb{R}^n} |\widetilde v(x,\omega)|^{\beta}\phi(x)\,\mathrm{d}x  
    \leq C_{1} \int_{\mathbb{R}^n}  \mathcal{R} \widetilde v (x,\omega)\cdot \nabla \phi(x)\,\mathrm{d}x + C_{2},
\end{equation}
where $C_1 = C_1(c_1,\beta)$ and $C_2(c_1,\beta,p)$ are constants that do not depend on $\omega$. We claim that $|\mathcal{R}^*\nabla \phi(x)|$ is bounded on $\R^n$ and there exists a constant $C_3$ depending only on $\phi$ (and hence on $c$ on $n$) such that for all $x$, 
\begin{equation}
\label{eq:bddR*phi}
    \left|\mathcal{R}^* \nabla \phi(x)\right| \le \frac{C_3}{|x|^{n+(\alpha-1)}}.%\frac{C_3}{|x|^{n+\delta+(\alpha-1)}}.
\end{equation}

With the uniform boundedness and the estimate above for $\mathcal{R}^*\nabla \phi$, and using the expression of $\phi$, we easily check
\begin{equation*}
\begin{aligned}
    C_4 := &\int_{\R^n} \left|\mathcal{R}^*\nabla \phi(x)\right|^{\beta'} \frac{1}{|\phi(x)|^{\beta'/\beta}}\,\dd x\\
    \le & \|\mathcal{R}^*\nabla \phi\|_{L^\infty}|B_1(0)| + C\int_{\{|x|\ge 1\}} \left(\frac{1}{|x|^{n+\alpha-1}} (\sqrt{1+|x|^2})^{\frac{n+\delta}\beta}\right)^{\beta'}\,\dd x\\
    < &\infty,
    \end{aligned}
\end{equation*}
as long as $\alpha > 1+\frac{\delta}{\beta}$ (i.e., $\delta < \beta(\alpha-1)$). Given $\alpha \in (1,2)$ and $\beta >1$, we can always choose $\delta>0$ so that this condition holds. Then applying H\"older's and then Young's inequalities, we get
\begin{equation*}
\begin{aligned}
    J := &\left|\int_{\mathbb{R}^n}  \mathcal{R} \widetilde v (x,\omega)\cdot \nabla \phi(x)\,\mathrm{d}x\right| = \left|\int_{\mathbb{R}^n} \widetilde v (x,\omega)\cdot \mathcal{R}^*\nabla \phi(x)\,\mathrm{d}x\right|\\
    \le & \left(\int_{\R^n} |\widetilde v(x)|^\beta \phi(x)\,\dd x\right)^{\frac1{\beta}}\left(\int_{\R^n} |\mathcal{R}^*\nabla \phi(x)|^{\beta'} \frac{1}{|\phi(x)|^{\beta'/\beta}}\,\dd x\right)^{\frac{1}{\beta'}}\\
    \le &\frac{1}{2C_1} \int_{\R^n}|\widetilde v(x,\omega)|^\beta \phi(x)\,\upd x + C_4\frac{1}{\beta'}(2C_1/\beta)^{\beta'/\beta}.
    \end{aligned}
\end{equation*}
Combining this estimate with \eqref{vLpE}, we get a constant $C_5=C_5(p,n,c_1,c_2,\alpha,\beta)$ such that
\begin{equation}
    \int_{\R^n} |\widetilde{v}(y,\omega)|^\beta \phi(y)\,\dd y \le C_5.
\end{equation}
This yields, for each $R>0$, a constant $C_R$ depending on $C_5$ and $R$ such that
\begin{equation*}
    \int_{B_R(0)} |\widetilde{v}(y,\omega)|^\beta \,\dd y \le C_R.
\end{equation*}
Finally, translating $\phi$ to $\phi(\cdot -x)$, we get the desired estimate \eqref{eq:unifbddv}.

\medskip

It remains to check the estimate \eqref{eq:bddR*phi}, and for this it suffices to consider $\mathcal{R}^*\partial_1 \phi$ as the other components are bounded in the same way. Let $\psi = \partial_1\phi$ and verify that $\psi$ is uniformly bounded and $\psi(x) \sim |x|^{-n-\delta-1}$ for large $|x|$. By the explicit form of the kernel of $\mathcal{R}$ in \eqref{eq:cRdef}, we see that for all $x\in \R^n$
\begin{equation*}
\begin{aligned}
    |\mathcal{R}^*\psi(x)| &= \left|\int_{\R^n} R(y-x)\psi(y)\,\dd y\right|\\ 
    &\le C\int_{\R^n} \frac{1}{|x-y|^{n+\alpha-2}}|\psi(y)|\,\dd y \le C(\|\psi\|_{L^\infty(\R^n)} + \|\psi\|_{L^1(\R^n)}).
    \end{aligned}
\end{equation*}
This verifies the uniform boundedness of $|\mathcal{R}^*\psi|$. Note that here and in the rest of the proof, the bounding constant $C$ will vary from line to line.

To prove the rate of $|\mathcal{R}^*\psi|$  for $|x|$ large, say $|x|\ge 4$, we first write the definition of $\mathcal{R}^*\psi$ as the sum of $I_1+I_2$ as follows:
\begin{equation*}
    I_1 = \int_{B_{|x|}(x/2)} R(y-x)\psi(y)\,\dd y,
\end{equation*}
and
\begin{equation*}
    I_2 = \int_{\R^n\setminus B_{|x|}(x/2)}  R(y-x)\psi(y)\,\dd y.
\end{equation*}

For $I_2$, we note that outside $B_{|x|}(x/2)$, 
\begin{equation*}
    \frac{|y|}{|y-x/2|}, \; \frac{|y-x|}{|y-x/2|} \quad \in \quad [\frac12, \frac32].
\end{equation*}
We then get
\begin{equation*}
    I_2 \le C\int_{\R^n\setminus B_{|x|}(x/2)} \frac{1}{|y-x/2|^{n+\alpha-2}} \frac{1}{|y-x/2|^{n+\delta+1}}\,\dd y \le C\frac{1}{|x|^{n+\delta+\alpha-1}}.
\end{equation*}

For $I_1$, we further break the integration domain into two regions. On $B_{|x|}(x/2)\setminus B_{|x|/2}(0)$, we have $|y|\ge |x|/2$, and hence $|\psi(y)|\le C|x|^{-n-\delta-1}$. We get
\begin{equation*}
\begin{aligned}
    \left| \int_{B_{|x|}(x/2)\setminus B_{|x|/2}(0)} R(y-x)\phi(y) \dd y\right| &\le C\frac{1}{|x|^{n+\delta+1}} \int_{B_{2|x|}(x)} |R(y-x)|\,\dd y\\
    &\le C\frac{1}{|x|^{n+\delta+1}} \int_{B_{2|x|}(0)} \frac{1}{|y|^{n-2+\alpha}}\,\dd y\\
    &\le C\frac{1}{|x|^{n+\delta+\alpha-1}}.
    \end{aligned}
\end{equation*}
For the remaining part of $I_1$, note that $R(y-x)$ and $\psi(y)$ are both smooth functions in this region (there is no singularity), and an integration by parts shows
\begin{equation*}
    \int_{B_{|x|/2}(0)} R(y-x)\psi(y)\,\dd y = \int_{\partial B_{|x|/2}(0)} n_1(y)R(y-x)\phi(y)\,\dd y - \int_{B_{|x|/2}(0)} \partial_1 R(y-x)\phi(y)\,\dd y.
\end{equation*}
For the integral on the sphere where by an abuse of notation $\dd y$ means the surface measure, we note $|n_1|\le 1$ and $|y-x|\ge |x|/2$. Hence on the sphere, $|R(y-x)|\le C|x|^{-n-\alpha+2}$ and $|\phi(y)|\le C|x|^{-n-\delta}$. It follows that
\begin{equation*}
    \left|\int_{\partial B_{|x|/2}(0)} n_1(y)R(y-x)\phi(y)\,\dd y\right| \le C\frac{1}{|x|^{n+\delta+\alpha-1}}.
\end{equation*}
For the other integral over $B_{|x|/2}(0)$, we note that $|y-x|>|x|/2$, and hence by the homogeneity of $\nabla R$ (of degree $-n-\alpha+1$), 
\begin{equation*}
    \left|\int_{B_{|x|/2}(0)} \partial_1 R(y-x)\phi(y)\,\dd y\right| \le C\frac{1}{|x|^{n+\alpha-1}} \int_{B_{|x|/2}(0)} \phi(y)\,\dd y \le C\frac{1}{|x|^{n+\alpha-1}}.
\end{equation*}
Combining all those estimates above, we verify \eqref{eq:bddR*phi} and complete the proof of the lemma.
\end{proof}

\end{document}
